%%
%%                  TEMPLATE for Math 436 project report
%%
\documentclass[11pt]{amsart}
%%% WARNING: Do NOT change the page size, fonts, or margins!  Penalties will apply.

% \usepackage{chemformula}
\usepackage[version=4]{mhchem}  % Used for chemical equations
\usepackage{graphicx}
\usepackage{amssymb,amsmath,amsthm}
\usepackage{placeins} %enables \FloatBarrier, which prevents figures and tables from going below it.
\usepackage{hyperref} %makes cross references into hyperlinks. 
\graphicspath{ {./images/} }

\begin{document}

\title{Go With the Flow: Optimal Sailboat Racing}
\author{Connor McBride, Ethan Palenske, Jackson Pond}

% Change the date to match the date you actually wrote this paper
\date{\today} % or use \today

\begin{abstract}
In this paper we will try to model the optimal path in the historical M\&M 100 Miler sailing race that takes place in Lake Michigan. Many forces impact optimal sailing, including wind direction, current, and waves, making this a challenging problem to model accurately. The primary questions we work to answer are what the simplest model is that rediscovers real sailing techniques, such as tacking, and how many parameters we can add into our model while still being able to solve it and getting reasonable results. We solve the classic Zermelo Navigation Problem assuming constant boat velocity and fixed current drift vector field. Further iterations involve relying on wind instead of a motor, taking into account the impact of the islands on wind and current speed, and including the impacts of being near Wisconsin on the wind.
\end{abstract}

\maketitle % this actually makes the title

%% First Section
\section{Background}

\par The primary question this paper seeks to answer is determining the sorts of models that accurately match the chemical reactions that yeast undergo within bread. Yeast interacts in complicated ways with the other ingredients in bread, and even with its own byproducts. We also want to determine what factors make failed models no longer accurate.

\par Every breadmaker has noticed that small changes in initial dough conditions can lead to wide variations in the quality of the baked bread. By having a good model for how yeast interacts with the other ingredients in bread, we will be able to better understand the breadmaking process and how to optimize initial conditions.

%% Second Section 
\section{Mathematical Representation}

\par In this section we go over the models we built in detail.

%% Third Section
\section{Solution}

\cite{modeling_yeast_spread}

%%Fourth Section
\section{Interpretation}


\subsection*{Appropriateness of models}


\subsection*{Potential improvements}


\subsection*{Impact}


%%%%%%%%%%%%%%%%%%%%%%%%%%%%%%%%%%%%%
%% Bibliography below
%%%%%%%%%%%%%%%%%%%%%%%%%%%%%%%%%%%%%
\FloatBarrier % Keep the figures from being put after the bibliography
\newpage
\bibliography{refs}
\bibliographystyle{alpha}

\end{document}
